\section{Multimodalidad nativa vs multimodalidad de puente}

\subsection{Comparación de las dos rutas}

El moderador introduce la discusión: multimodalidad nativa (convierte directamente los pixel en token y, desde el inicio, se entrena con una mezcla de imagen y texto) vs multimodalidad de puente (primero se entrena el LLM y después se conecta un ViT + adapter).

\begin{importantbox}{Xue Fuzhao (薛福昭): Native es más clean, pero el Encoder tiene ventajas de ingeniería}
\textbf{Ventajas de la multimodalidad de puente (Encoder-based)}:
\begin{itemize}
  \item El Encoder y el Decoder se pueden entrenar por separado, lo que facilita la optimización
  \item Se puede comprimir de forma flexible la cantidad de Token, con mejor eficiencia de cómputo global
\end{itemize}
\textbf{Desventajas de la multimodalidad de puente}:
\begin{itemize}
  \item El objetivo de entrenamiento del Encoder (¿CLIP? ¿Detection? ¿Reconstruction?) no puede garantizar que sea completamente sin pérdida para las etapas posteriores
  \item Alta complejidad de Infra: la posición irregular de Image provoca un workload desbalanceado durante el Sequence Parallelism, lo que exige un reshuffle/recompute complejo y daña mucho el MFU
\end{itemize}
\textbf{Ventajas de Native}: el pixel token y el text token no tienen diferencia, y a nivel de Infra casi no es necesario modificar la infraestructura de un LM puro.
\end{importantbox}

\textbf{Liu Ziwei (刘子伟) (prefiere Native)}: las tres grandes ventajas de Native:
\begin{enumerate}
  \item \textbf{Consideración global de Data} — no hace falta entrenar primero el ViT y luego conectar el LM; la proporción de datos es más unificada
  \item \textbf{Emerging Ability} — desde el inicio hubo Interleaved Data, y el Early Fusion puede producir capacidades emergentes (se reporta que Gemini 3 volvió a entrenar la versión nativa)
  \item \textbf{Infra Clean} — se puede reutilizar directamente el Infra del LM, y el despliegue es más simple
\end{enumerate}
Pero por ahora falta un Codebase/Checkpoint de Native Multimodal de código abierto que todos puedan usar.

\textbf{Xie Tianbao (谢天宝)}: la Native Multimodal se discute desde Gemini 1 en 2,023. En esencia es un problema de madurez de ingeniería — la industria ya tiene las ideas, pero la capacidad de ingeniería todavía no las alcanza del todo, y sigue aprovechando los beneficios del esquema actual. \textbf{Al final, cuál esquema se queda depende de la mantenibilidad y la controlabilidad} (por analogía con la ingeniería de software).

\textbf{Cao Yu (曹玉) (perspectiva de datos)}: el desarrollo de los modelos multimodales \textbf{sigue la disponibilidad de los datos}. La multimodalidad nativa tiene mejor afinidad con los datos — para el posentrenamiento, cuanto más fuerte sea la capacidad base del base model, mejor, y la arquitectura native tiene una ventaja natural en este aspecto. Al mismo tiempo plantea una pregunta sin resolver: ¿qué modalidades debería incluir el \textbf{lado de salida} de la multimodalidad nativa? En el lado de entrada (text/vision/audio) no hay controversia, pero en el lado de salida cada quien opina distinto — los seres humanos mismos no tienen una capacidad de salida vision native.

\subsection{¿Por qué se separan LLM y VLM? ¿Cuándo se unificarán?}

El moderador señala que muchas empresas nacionales todavía publican el LLM y el VLM por separado, mientras que en el extranjero (GPT, Gemini) se tiende a modelos unificados.

\textbf{Xue Fuzhao}: los modelos de código cerrado son una caja negra; puede que internamente sea simplemente un enrutamiento entre dos modelos. Si se despliega con una estructura Encoder + Decoder, lo mejor es separarlos en dos chips — una query solo de texto (text-only) no necesita la memory del vision encoder. Ponerlos juntos complica el problema, pero a nivel de Serving es resoluble.

\textbf{Xie Tianbao}: la razón por la que se separan LLM/VLM es la misma por la que se separa el desarrollo de Coding/Math/Tool Use — desarrollar por ramas y después fusionar es una práctica común en ingeniería. Por ahora se está en un estado de «versión intermedia».

\textbf{Cao Yu}: la separación de VLM y LLM es temporal. Para una entidad operante con verdadera inteligencia, la capacidad de comprender el mundo \textbf{definitivamente no puede ser de una sola modalidad}. La separación actual se debe a que el Transformer modela discrete signal de forma mucho más madura que continuous signal. A largo plazo se superarán los problemas de ingeniería/organización/tecnología y se avanzará hacia una multimodalidad de entrada nativa.

\textbf{Liu Ziwei}: la separación a corto plazo tiene factores históricos (el Benchmark se inclina hacia la inteligencia textual, y Vision no influye mucho en la clasificación). Pero desde la perspectiva de Physical AI — que necesita decisiones de largo alcance (long-horizon action/decision) y fusión multimodal —, estos escenarios \textbf{deben unificarse}. Los seres humanos «primero tienen capacidad motora y después lenguaje», la IA es al revés, y al final también tendrá que resolver esas tareas de base que necesitan sensory input.

\subsection{Resumen de la sección}
Native vs Encoder-based cada uno tiene ventajas y desventajas; a corto plazo coexisten y a largo plazo la tendencia es hacia Native. La separación de LLM/VLM es un estado temporal de ingeniería/organización; en la era de Physical AI se unificarán necesariamente.

%% ===== Datos sintéticos y Mid-training =====
